Table~\ref{tab:robust} in Appendix~\ref{app:exp} varies one choice at a time, with 30 seeds for the certified search
and 20 for the baselines.

\paragraph{Budget.} A smaller budget share makes pilots shallower and certificates slower. With the monotone class
the final rate reached \num{rob:budget05:certM:final} of the optimum with $\lrel=0.05$ and
\num{rob:budget20:certM:final} with $\lrel=0.2$, against \num{main:certM:final} with $\lrel=0.1$. No pilot exceeded
its budget in either variant. With $\lrel=0.2$, \num{rob:budget20:harmall} adoptions lowered profit, one for each
class in the same market and seed, after the same interval had missed.

\paragraph{Confidence intervals.} With Hoeffding intervals instead of relative-entropy intervals, the monotone class
reached \num{rob:hoeffding:certM:final} of the optimum. The loss is concentrated in the markets with two or three
competing quotes, where take-up is low and Hoeffding intervals are several times wider; with one competing quote
the two kinds of interval give almost the same result.

\paragraph{History.} With 50{,}000 quotes of history at $b_0$ before the first pilot, the monotone class reached
\num{rob:hist:certM:final} of the optimum. History narrows the incumbent's interval and so allows deeper first
pilots, but it does not remove the need for pilots (Theorem~\ref{thm:nocert}).

\paragraph{A wrong slope bound.} With a slope bound half the true maximum slope, the class $\Meta$ no longer contains
the take-up curve: \num{rob:etahalf:certeta:pilotover} pilots exceeded their budget, in
\num{rob:etahalf:certeta:runsover} of runs, and \num{rob:etahalf:certeta:over} of quotes were over budget. The
overruns can be large: the worst pilot lost \num{rob:etahalf:certeta:pilotlossmax} of the incumbent's profit per
quote, and under the convex class of the main setting the worst lost \num{main:certC:pilotlossmax}. No adoption
lowered profit (\num{rob:etahalf:certeta:harm} harmful adoptions): a wrong class can make pilots much costlier than
budgeted, but adoptions do not use it (Theorem~\ref{thm:anytime}(i)).

\paragraph{The own-data prior.} The failure of the sampler that starts from the own-data model comes from the
precision of its prior, not from its starting point alone. With prior standard deviations ten times larger its
final rate reached \num{rob:tsown10:tsown:final} of the optimum and its average profit
\num{rob:tsown10:tsown:avg}, and with a hundred times larger \num{rob:tsown100:tsown:final} and
\num{rob:tsown100:tsown:avg}, against \num{main:tsown:final} and \num{main:tsown:avg} in the main setting. The main
setting gives the slope the standard error that the lender's files report. That standard error is small because the
files hold many applications, but all of them come after the borrower chose the lender, so the precision is about
the wrong quantity. Nor is the failure due to the level of the prior, which is the take-up of approved offers
(about 97\%) and far above take-up at the quote stage: with the level set to the true take-up at $b_0$ and the
same flat slope, the final rate reached \num{rob:tsownlevel:tsown:final} of the optimum. The sampler raises its
rate, its quotes then sit in a narrow range of high rates where the data say little about the slope, and
500{,}000 quotes do not overturn a prior that tight.

\paragraph{Arms scaled to the market.} The baselines' arms cover three quarters of a point on each side of $b_0$
whatever the margin, so in the markets with thin margins most arms are far from the optimum. With arms at
multiples of one eighth of the optimal margin, from one eighth to three times the margin, which uses knowledge a
lender does not have, KL-UCB earned \num{rob:armsmargin:klucb:avgone} of the optimum on average and Thompson
sampling with a vague prior \num{rob:armsmargin:ts:avgone}; \num{rob:armsmargin:klucb:anyover} and
\num{rob:armsmargin:ts:anyover} of their runs still contained a quote over budget
(\num{rob:armsmargin:klucb:over} and \num{rob:armsmargin:ts:over} of quotes).

\paragraph{Design parameters of the search.} The pilot cap, the checkpoint schedule, the rate increment and the
error level were set once, after preliminary runs on a single market. Varying each one
(Table~\ref{tab:design}, \num{rob:design:variants} variants, \num{rob:design:runs} runs) moves the final profit
of the monotone class between \num{rob:design:certM:lo} and \num{rob:design:certM:hi} of the optimum and that of
the slope-bound class between \num{rob:design:certeta:lo} and \num{rob:design:certeta:hi}. The guarantees do not
depend on these parameters: across all variants \num{rob:design:pilotover} pilots exceeded their budget (exact 95\%
upper limit \num{rob:design:pilotover:cp:hi} per seed, from \num{rob:design:seeds} seeds), and
\num{rob:design:harm} adoptions lowered profit, each in a run where an interval had missed.
